\paragraph{Forward projection to Watson constants.}

Once $\tau'$ has been constructed, one first defines the Cartesian
quartic matrix

\begin{equation}
T_{ij}=\frac{\tau'_{ij}}{4},
\label{eq:T_from_tauprime}
\end{equation}

then forms the cylindrical quartic combinations

\begin{align}
T_{400} &= \frac{3T_{aa}+3T_{bb}+2T_{ab}}{8}, \\
T_{220} &= T_{ac}+T_{bc}-2T_{400}, \\
T_{040} &= T_{cc}-T_{220}-T_{400}, \\
T_{202} &= \frac{T_{aa}-T_{bb}}{4}, \\
T_{022} &= \frac{T_{ac}-T_{bc}}{2}-T_{202}, \\
T_{004} &= \frac{T_{aa}+T_{bb}-2T_{ab}}{16},
\label{eq:cylindrical_quartic}
\end{align}

and finally evaluates the quartic constants in Watson's $A$ reduction,

\begin{align}
\Delta_J   &= -T_{400}-2T_{004}, \\
\Delta_{JK}&= -T_{220}+12T_{004}, \\
\Delta_K   &= -T_{040}-10T_{004}, \\
\delta_J   &= -T_{202}, \\
\delta_K   &= -T_{022}-4\Sigma T_{004},
\label{eq:watson_A_quartic}
\end{align}

with

\begin{equation}
\Sigma=\frac{2A-B-C}{B-C}.
\end{equation}

Equivalently, the $S$-reduced constants are obtained as

\begin{align}
D_J   &= -T_{400}+\frac12 T_{022}\Sigma_1, \\
D_{JK}&= -T_{220}-3T_{022}\Sigma_1, \\
D_K   &= -T_{040}+\frac52 T_{022}\Sigma_1, \\
d_1   &= T_{202}, \\
d_2   &= T_{004}+\frac14 T_{022}\Sigma_1,
\label{eq:watson_S_quartic}
\end{align}

with

\begin{equation}
\Sigma_1=\frac{B-C}{2A-B-C}=\Sigma^{-1}.
\end{equation}

The forward projection from the tensorial description to Watson's
quartic constants is therefore unique once the reduced quartic tensor
$\tau'$ has been specified.

\paragraph{Inverse reconstruction and affine gauge family.}

The inverse problem has a different structure.
The spectroscopic constants in a fixed Watson reduction provide only
five independent quartic parameters, whereas the underlying pair--pair
quartic tensor contains six independent coordinates.
Consequently, the linear map

\begin{equation}
\boldsymbol\tau^{(4)} \longmapsto \mathbf D^{(4)}
\end{equation}

has rank five and therefore possesses a one-dimensional null space.

In the $A$ reduction, if

\begin{equation}
\Sigma=\frac{2A-B-C}{B-C},
\end{equation}

a null vector of the projection operator can be written as

\begin{equation}
\mathbf n =
\left(
0,\,
0,\,
0,\,
1,\,
\frac{\Sigma-1}{2},\,
-\frac{\Sigma+1}{2}
\right)^T.
\label{eq:quartic_null_vector}
\end{equation}

In terms of the reduced tensor $\tau'$, the corresponding null matrix is

\begin{equation}
N'=
\begin{pmatrix}
0 & 1/2 & (\Sigma-1)/4 \\
1/2 & 0 & -(\Sigma+1)/4 \\
(\Sigma-1)/4 & -(\Sigma+1)/4 & 0
\end{pmatrix}.
\label{eq:quartic_null_matrix}
\end{equation}

Therefore, if $\tau'$ is any quartic tensor consistent with a given set
of Watson constants, then

\begin{equation}
\tau' \;\sim\; \tau' + \eta N'
\label{eq:quartic_affine_family}
\end{equation}

for arbitrary scalar $\eta$ produces exactly the same spectroscopic
quartic constants.

This is the key structural point: the one-dimensional null space does
not represent a physical constraint on quartic centrifugal distortion.
Rather, it defines an affine gauge family of tensor representatives
associated with a fixed set of Watson constants. The inverse
reconstruction of the tensor from spectroscopic constants therefore
requires the selection of a gauge within this affine family. In the
present framework the Moore--Penrose pseudoinverse provides a natural
and numerically stable gauge fixing, selecting one distinguished
representative within this family.
